\begin{table}
\begin{threeparttable}[t]
\centering
\caption{Standardized Effect Sizes}
\label{tab:sde}
\small
\begin{tabular}{lcccccc}
\toprule
Outcome & $\hat{\beta}$ & SE & SD($Y$) & SDE & SE(SDE) & Classification \\
\midrule
\multicolumn{7}{l}{\textit{Panel A: Pooled}} \\[3pt]
Competitive share & 0.2562 & 0.0943 & 0.2032 & 0.280 & 0.103 & Large positive \\
Single-bidder rate & -0.0665 & 0.0143 & 0.4166 & -0.160 & 0.034 & Large negative \\
\midrule
\multicolumn{7}{l}{\textit{Panel B: Heterogeneous (department size splits)}} \\[3pt]
Comp.\ share (large depts) & 0.2129 & 0.0579 & 0.1612 & 0.268 & 0.073 & Large positive \\
Comp.\ share (small depts) & 0.2387 & 0.1713 & 0.2344 & 0.247 & 0.177 & Large positive \\
\bottomrule
\end{tabular}
\begin{tablenotes}[flushleft]\small
\item \textit{Notes:} \textbf{Country:} Colombia. \textbf{Research question:} Does transitioning government procurement from an informational platform to a fully transactional e-procurement system increase competitive bidding in public contracts? \textbf{Policy mechanism:} Colombia's SECOP~II platform (staggered rollout 2015--2022) replaced SECOP~I, a publication-only portal, with a fully transactional interface where firms submit bids, upload documents, and receive awards entirely online, reducing the transaction costs of participating in competitive procurement processes. \textbf{Outcome definition:} Competitive share is the fraction of a department's quarterly contracts awarded through competitive modalities (licitaci\'{o}n p\'{u}blica, selecci\'{o}n abreviada, m\'{i}nima cuant\'{i}a, concurso de m\'{e}ritos) versus direct contracting. \textbf{Treatment:} Continuous --- SECOP~II share is the fraction of department-quarter contracts processed on the transactional platform (varies 0 to 1). \textbf{Data:} Colombia Compra Eficiente SECOP Integrado, 2015--2021, department-quarter panel, 1,130 department-quarter observations across 38 departments. \textbf{Method:} Difference-in-differences with department and year-quarter fixed effects; standard errors clustered at department level. \textbf{Sample:} All departments with at least 4 quarters of procurement data. SDE $= \hat{\beta} \times \text{SD}(X) / \text{SD}(Y)$ for continuous treatment. Classification refers to magnitude, not statistical significance: Large ($|$SDE$| > 0.15$), Moderate ($0.05$--$0.15$), Small ($0.005$--$0.05$), Null ($< 0.005$).
\end{tablenotes}
\end{threeparttable}
\end{table}
